%%%%% Section 5. Learning the Bundle from Data %%%%%
\section{Computing Facet-induced OT-FB}\label{sec:computation}

%%% Subsection 5.1 Objective function %%%
\subsection{Objective Function and Uniqueness}

We estimate the facet-induced OT-FB by fitting the model distribution
$q_\Phi(x)$, defined by the embedding map $\Phi$, to the content data
distribution $p_\text{data}(x)$:
\begin{align*}
    p_{\text{data}}(x) & = \sum_{n=1}^N p(n)\,
    \mathcal{N}(x \mid \mathbf{x}_n, \beta^{-1}I), \quad
    p(n) = \frac{1}{K N_{k(n)}},                                \\
    q_\Phi(x)          & = \iint_{\mathcal{U}\times\mathcal{V}}
    \mathcal{N}(x \mid \Phi(u, v), \beta^{-1}I)\, q(u)\, q(v)\, du\, dv,
\end{align*}
where $\beta^{-1}$ is the observation noise variance, $k(n)$ the genre
of work $n$, $N_k$ the number of works in genre $k$, and $q(u)$,
$q(v)$ the priors on the two spaces. We estimate $\Phi$ together with
the work coordinates $U = (u_1, \dots, u_N)$ in $\mathcal{U}$ and
$V = (v_1, \dots, v_N)$ in $\mathcal{V}$.

The objective aligns the $K$ genre distributions by minimizing the
2-Wasserstein distance---a multi-marginal optimal transport
problem~\citep{Pass20151771}.
\begin{equation}\label{eq:main_objective}
    \begin{aligned}
        W^2_2(p_\text{data}, q_\Phi)
               & = \inf_{\gamma\in\Pi(p_{\text{data}},\, q_\Phi)}
        \iint \|x - x'\|^2\, \gamma(x, x')\, dx\, dx',            \\
        \Phi^* & = \argmin_{\Phi} W^2_2(p_\text{data}, q_\Phi),
    \end{aligned}
\end{equation}
where $\gamma \in \Pi$ is a transport plan meeting the two marginal
constraints $p_{\text{data}}(x) = \int \gamma(x,x')\,dx'$
($p \rightarrow q$) and $q_\Phi(x') = \int \gamma(x,x')\,dx$
($q \rightarrow p$). Because $\Phi$ shares one connotative coordinate
$v$ across all genres, a single minimization fits every genre manifold
and optimizes the correspondence among them at once. The
correspondence is unique in the following sense.

\begin{proposition}\label{prop:OT-FB}
    Let each genre distribution $p_k$ be smoothed by the noise model
    $\mathcal{N}(x \mid \mathbf{x}_n, \beta^{-1}I)$, so that it is
    absolutely continuous with finite second moment. Then:
    (a) the multi-marginal optimal transport problem over
    $p_1, \dots, p_K$ with pairwise squared-distance cost has a unique
    solution, given by Monge maps; the OT-FB is therefore determined by
    the collection data.
    (b) The map $\Phi$ realizing it is unique up to relabeling of the
    coordinates of $\mathcal{U}$ and $\mathcal{V}$, which leaves the
    arrangement of works unchanged.
\end{proposition}

\begin{proof}[Proof sketch]
    (a) Each $p_k$ is a finite Gaussian mixture, hence absolutely
    continuous with finite second moment. For such marginals, the
    multi-marginal problem with pairwise squared-distance cost has a
    unique solution, given by Monge maps normalized at one
    marginal~\citep{GangboSwiech199823}; for $K=2$ this is
    Brenier's theorem~\citeyearpar{Brenier1991375}. Solutions depend on the
    cost, and \citet{Pass20151771} surveys the general case. The OT-FB
    inherits this uniqueness.
    (b) Let $g_\mathcal{U}$ and $g_\mathcal{V}$ be measure-preserving
    maps of $\mathcal{U}$ and $\mathcal{V}$, leaving $q(u)$ and $q(v)$
    invariant. Then $\Phi \circ (g_\mathcal{U} \times g_\mathcal{V})$
    gives the same $q_\Phi$, the same transport cost, and the same
    correspondence. Conversely, any two maps realizing the
    correspondence of (a) differ only by such a reparametrization.
    \hfill\openbox
\end{proof}

This is the sense of \emph{essentially unique} used in the main text:
the collection data fix the arrangement of works, not the labels of
the axes.

%%% Subsection 5.2 Computing OT-FB from data
\subsection{Optimization}

We minimize \eqref{eq:main_objective} over $\Phi$ and the coordinates
$(u_n, v_n)$ of every work $n$, subject to the two marginal
constraints. We enforce the constraints alternately, each step under
one-sided (unbalanced) transport.

\paragraph{Transport optimization $q_\Phi \rightarrow p_\text{data}$.}
With $(U, V)$ fixed and the constraint $q \rightarrow p$ satisfied,
the minimizing map is
\begin{equation}
    \Phi(u, v \mid U, V) =
    \frac{\sum_k \mathcal{N}(u \mid u_k, \alpha^{-1} I)\, \varphi_k(v)}
    {\sum_{k'} \mathcal{N}(u \mid u_{k'}, \alpha^{-1} I)},
    \label{eq:kernel_smoothing}
\end{equation}
where $u_k$ is the denotation coordinate shared by the works of genre
$k$ (that is, $u_n \equiv u_{k(n)}$ under hard genre labels), and
$\varphi_k$ is the manifold model of genre $k$:
\begin{equation}\label{eq:varphi}
    \varphi_k(v) =
    \frac{\sum_{n:\rho_{kn}=1} \mathcal{N}(v \mid v_n, \alpha^{-1} I)\,
        \mathbf{x}_n}
    {\sum_{n':\rho_{kn'}=1} \mathcal{N}(v \mid v_{n'}, \alpha^{-1} I)}.
\end{equation}
Here $\rho_{kn} = 1$ when work $n$ belongs to genre $k$ and $0$
otherwise, and $\alpha$ controls the kernel width. Equation
\eqref{eq:kernel_smoothing} is \textit{hierarchical kernel
    smoothing}---the algorithm of the SOM of
SOMs~\citep{Furukawa2009463}---obtained here not as a modeling
assumption but as a consequence of optimal transport.

\paragraph{Transport optimization $p_\text{data} \rightarrow q_\Phi$.}
Substituting \eqref{eq:kernel_smoothing} into
\eqref{eq:main_objective} under the constraint $p \rightarrow q$
yields an objective over $(U, V)$:
\begin{equation}\label{eq:E-UV}
    E[U, V] = \frac{\beta}{2}\sum_{n=1}^N p(n)
    \iint p(u, v \mid n)\,
    \bigl\|\mathbf{x}_n - \Phi(u, v \mid U, V)\bigr\|^2
    \, du\, dv
    + R(U, V),
\end{equation}
where $p(u, v \mid n)=\mathcal{N}(u\mid u_n,\alpha^{-1} I)\,\mathcal{N}(v\mid v_n,
    \alpha^{-1} I)$ is the distribution the transport plan assigns to
work $n$, and $R$ is the regularization term
\begin{equation*}
    R(U, V) \equiv
    - \lambda_1 \sum_{n=1}^N p(n)\,
    \bigl(\log q(u_n) + \log q(v_n)\bigr)
    - \lambda_2 \left(
    \int q(u)\log p(u)\, du
    + \sum_k \int q(v)\log p_k(v)\, dv
    \right).
\end{equation*}
The first term is the log prior, holding each work inside $\mathcal{U}$ and
$\mathcal{V}$. The second is an entropy term, spreading the works over both:
$p(u)$ and $p_k(v)$ are the densities the works make there.

\paragraph{Final objective.}
The Laplace approximation gives
\begin{equation}\label{eq:reduced_objective}
    E[U, V] \simeq \sum_{n=1}^N p(n) \left\{
    \frac{\beta}{2}\bigl\| \mathbf{x}_n - \Phi(u_n, v_n \mid U, V)
    \bigr\|^2
    \right\}+ R(U,V).
\end{equation}
A single objective over $(U, V)$ remains, which we minimize by
gradient descent.



%%% Subsection 5.3 Multiple facets %%%
\subsection{Multiple Facets}

Two facets make a class a pair: the $i$-th focus of one and the $j$-th
of the other. The denotation coordinate splits, $u_{ij} = (u^1_i,
    u^2_j)$, and so does the kernel of \eqref{eq:kernel_smoothing}:
\begin{equation*}
    \mathcal{N}(u \mid u_{ij}, \alpha^{-1} I) =
    \mathcal{N}(u^1 \mid u^1_i, \alpha^{-1} I)\,
    \mathcal{N}(u^2 \mid u^2_j, \alpha^{-1} I).
\end{equation*}
Smoothing runs along each facet in turn, following the tensor
SOM~\citep{Iwasaki2016107}, and $\varphi_{ij}$ is fitted to each cross
class. More facets follow the same way.